%%%%%%%%%%%%%%%%%%%%%%%%%%%%%%%%%%%%%%%%%
% University/School Laboratory Report
% LaTeX Template
% Version 4.0 (March 21, 2022)
%
% This template originates from:
% https://www.LaTeXTemplates.com
%
% Authors:
% Vel (vel@latextemplates.com)
% Linux and Unix Users Group at Virginia Tech Wiki
%
% License:
% CC BY-NC-SA 4.0 (https://creativecommons.org/licenses/by-nc-sa/4.0/)
%
%%%%%%%%%%%%%%%%%%%%%%%%%%%%%%%%%%%%%%%%%

%----------------------------------------------------------------------------------------
%	PACKAGES AND DOCUMENT CONFIGURATIONS
%----------------------------------------------------------------------------------------

\documentclass[
	letterpaper, % Paper size, specify a4paper (A4) or letterpaper (US letter)
	11pt, % Default font size, specify 10pt, 11pt or 12pt
]{CSUniSchoolLabReport}

\usepackage{subcaption}
\usepackage{booktabs}
\geometry{top=2cm} % Override top margin to reduce space before title

%----------------------------------------------------------------------------------------
%	REPORT INFORMATION
%----------------------------------------------------------------------------------------

\title{Assignment 2} % Report title

\author{Yiran \textsc{Liu}} % Author name(s), add additional authors like: '\& James \textsc{Smith}'

\date{\today} % Date of the report

%----------------------------------------------------------------------------------------

\begin{document}

\maketitle % Insert the title, author and date using the information specified above

\begin{center}
	\begin{tabular}{l r}
      ECE 108: Discrete Math \& Logic 1 \\
      Student Number: 21184901
	\end{tabular}
\end{center}
%----------------------------------------------------------------------------------------
%	COMPUTATIONS
%----------------------------------------------------------------------------------------

\subsubsection*{Question 1}
\textit{Prove, for every $n\in\mathbb{Z}^+$, that $7^n - 2^n$ is divisible by $5$:} \\

\textbf{Solution:} We prove by induction on $n$. Suppose $7^n - 2^n = 5i$ where $i \in \mathbf{Z^+}$:

\begin{align*}
	7^{n+1} - 2^{n+1} & = 7^n \times 7 - 2^n \times 2 \\
	& = 7 \times (7^n - 2^n \times \frac{2}{7}) \\
	& = 7 \times (7^n - 2^n + \frac{5}{7} \times 2^n) \\
	& = 7 \times (7^n - 2^n) + 5 \times 2^n \\
	& = 7 \times (5i) + 5 \times 2^n \\
	& = 5 (7i + 2^n)
\end{align*}

This proves that $7^n - 2^n \pmod{5} = 0 \implies 7^{n+1} - 2^{n+1} \pmod{5} = 0$

We observe that for $n=1$, $7^1 - 2^1 = 5$ is divisible by 5. Therefore, we can induce that $7^n - 2^n$ is divisible by $5$ for $\forall n\in\mathbb{Z}^+$.

\pagebreak

\subsubsection*{Question 2}
\textit{Prove, for every $n\in\mathbb{Z}^+$, where $n = \left(n_{k-1}\,n_{k-2}\,\ldots\,n_0\right)_{10}$, where each $n_i \in\{0,\ldots, 9\}$ is a digit, that:}
\[
n\text{ is divisible by } 3 \implies \sum_{i=0}^{k-1} n_i \text{ is divisible by } 3
\]
\textbf{Solution:} For the sake of simplicity, we only consider one way to spell the digits of a number by never putting 0 in the front.
We define \[
	\mathrm{digisum}(n) = \sum_{i=0}^{k-1} n_i.
\]

Observe that for a number $n$ with digits $(n_{k-1}, n_{k-2}, \ldots, n_0)$, if we increase $n$ by $3$:
\begin{enumerate}
	\item[\textcircled{1}] For $n_0 < 7$: $n_0 + 3 < 10$, $n_0$ increases by $3$ and all other digits stay the same,
		$\mathrm{digisum}(n+3) = \mathrm{digisum}(n) + 3$
	\item[\textcircled{2}] For $n_0 \geq 7$ and $n_1 \neq 9$: $n_0' = n_0 + 3 - 10 = n_0 - 7$, $n_1$ increases by $1$,
		$\mathrm{digisum}(n+3) = \mathrm{digisum}(n) - 6$
	\item[\textcircled{3}] For $n_0 \geq 7$ and $n_1 = 9$: $n_0' = n_0 + 3 - 10 = n_0 - 7$, $n_1' = 0$, and:
		\begin{itemize}
			\item For the first $a$ digits where $n_1, \ldots, n_a = 9$: $n_i' = n_i - 9 = 0$
			\item For the very first digit that $n_i \neq 9$: $n_i' = n_i + 1$
		\end{itemize}
		\[ \mathrm{digisum}(n+3) = \mathrm{digisum}(n) - 7 - 9a + 1 = \mathrm{digisum}(n) - 3(2+3a) \]
		where $a \in \mathbb{Z}^+$
\end{enumerate}

For each case in \textcircled{1}\textcircled{2}\textcircled{3}: $\mathrm{digisum}(n+3) \pmod 3 = \mathrm{digisum}(n) \pmod 3$.

For the base case ($n=3$), $\text{digisum(3)=3}$. We can therefore induce that $\forall n \in \mathbf{Z^+} \mid n \pmod{3} = 0$, $\text{digisum}(n) \pmod{3} = 0$ 

\pagebreak

\subsubsection*{Question 3}
\begin{itemize}
	\item Suppose $S = \emptyset$; i.e., $S$ is the empty set. Then is the following proposition true? \begin{quote}If $x\in S$, then $x$ is the Pope.\end{quote}
	\textbf{Solution: } Yes. Since $x \in \emptyset$ is always false, the proposition $\text{false} \implies \text{false}$ is vacuously \textbf{true}. 
	\item With the restriction $\mathcal{X} \not= \emptyset$ no longer enforced and with $\mathbb{N}$ as the universal set, what is $\bigcap \mathcal{X}$ when $\mathcal{X} = \emptyset$?

	\textbf{Solution: } $\bigcap \mathcal{X} = \{x \in \mathbb{N} \mid \forall A \in \mathcal{X}, x \in A\}$. But since $\mathcal{X} = \emptyset$ there is no such $A \in \mathcal{X}$. Again, the proposition $\text{false} \implies \text{false}$ is \textbf{true}. Therefore, $\bigcap \mathcal{X} = \mathbb{N}$
\end{itemize}

\pagebreak

\subsubsection*{Question 4}
\textit{Prove: $A\cap \left(A \cup B\right) = A \cup \left(A \cap B\right) = A$} \\

\textbf{Solution: } The first half of the equation can be proven directly using distributive law: \[
	A\cap (A\cup B) = (A \cap A) \cup (A \cap B) = A \cup (A \cap B)
	\]
As for the second half, by definition: \[
	x \in A \cup (A \cap B)
	\iff (x \in A) \lor ((x \in A) \land (x \in B)) \iff (x \in A)
	\]
Therefore, $A = A \cup (A \cap B)$

\pagebreak

\subsubsection*{Question 5}
\textit{Prove: $\left(A\subseteq B\right) \iff \left(A\cap B = A\right) \iff \left(A\cup B = B\right) \iff \left(A\cap\overline{B} = \emptyset\right) \iff \left(\overline{A} \cup B = \mathcal{U}\right)\iff \left(\overline{B}\subseteq \overline{A}\right) \iff \left(A\setminus B = \emptyset\right)$.}

\textbf{Solution: } We can write the equivalence of these propositions according to the definition of sets: 
\[
\left\{
\begin{aligned}
	A \subseteq B &\iff (x \in A \implies x \in B) \\
	A \cap B = A &\iff (x \in A \land x \in B \iff x \in A) \\ 
	A \cup B = B &\iff (x \in A \lor x \in B \iff x \in B) \\
	A \cap \overline{B} = \emptyset &\iff (x \in A \land \lnot(x \in B) \iff x \in \emptyset) \\
	\overline{A} \cup B = \mathcal{U} &\iff (\lnot(x \in A) \lor x \in B \iff x \in \mathcal{U}) \\
	\overline{B} \subseteq \overline{A} &\iff (\lnot(x \in B) \implies \lnot(x \in A)) \\
	A \setminus B = \emptyset &\iff (x \in A \land \lnot(x \in B) \iff x \in \emptyset)
\end{aligned}
\right.
\]

Let P and Q denote the propositions $x \in A$ and $x \in B$, respectively, our statement becomes theories proved in the textbook:

\[
\begin{aligned}
	  &(\text{P} \implies \text{Q}) \\
	\iff{} &((\text{P} \land \text{Q}) \iff \text{P}) \\
	\iff{} &((\text{P} \lor \text{Q}) \iff \text{Q}) \\
	\iff{} &((\text{P} \land \lnot \text{Q}) = \text{false}) \\
	\iff{} &((\lnot \text{P} \land \text{Q}) = \text{true}) \\
	\iff{} &(\lnot \text{Q} \implies \lnot \text{P}) \\
	\iff{} &((\text{P} \land \lnot \text{Q}) = \text{false})
\end{aligned}
\]

\pagebreak

\subsubsection*{Question 6}
Suppose $A, B, C$ are sets. First clearly state true or false, then prove or disprove:
\begin{itemize}
	\item $A\cap B = \emptyset \iff A = A\setminus B$. \textbf{True.}

	\textbf{Solution: } Since $A \setminus B = \{x \in A \mid x \notin B\}$:
	\begin{itemize}
		\item If $A \cap B = \emptyset$, then for any $x \in A$, $x \notin B$, so $x \in A \setminus B$. Also $A \setminus B \subseteq A$. Therefore $A = A \setminus B$.
		\item If $A = A \setminus B$, every $x \in A$ satisfies $x \notin B$, so $A \cap B = \emptyset$.
	\end{itemize}

	\item $A\setminus\left(B\setminus C\right) = A\cap \left(\overline{B} \cup C\right)$. \textbf{True.}

	\textbf{Solution: }
	\[
	\begin{aligned}
		x \in A\setminus(B\setminus C)
		&\iff (x \in A) \land \lnot\left((x \in B) \land (x \notin C)\right) \\
		&\iff (x \in A) \land \left((x \notin B) \lor (x \in C)\right) \\
		&\iff x \in A \cap \left(\overline{B} \cup C\right)
	\end{aligned}
	\]

	\item $A \subseteq B \implies A\times C \subseteq B\times C$. \textbf{True.}

	\textbf{Solution: } Suppose $A \subseteq B$. Let $(x, y) \in A \times C$; then $x \in A$ and $y \in C$.
	Since $A \subseteq B$, $x \in B$, so $(x, y) \in B \times C$.

	\item $A\times C \subseteq B\times C \implies A\subseteq B$. \textbf{False.}

	\textbf{Solution: } Counterexample: $A = \{1\}$, $B = \{2\}$, $C = \emptyset$.
	Then $A \times C = \emptyset \subseteq \emptyset = B \times C$, but $1 \in A$ and $1 \notin B$, so $A \not\subseteq B$.
\end{itemize}

\pagebreak

\subsubsection*{Question 7}

\textit{Suppose $\mathcal{X} = \{[a, \infty) \mid a < 0\}$. Prove that $\bigcup \mathcal{X} = \mathbb{R}$.} \\

\textbf{Solution: } Since for $\forall x \in \mathbf(R)$, we can find $\exists a \in \mathbf(R), a < 0$ such that $y < x$. We can find at least one $\exists a \in \mathbf{R}, a < 0$ such that $x \in [a, \infty)$. 
This means that $x \in \mathbf{R} \implies (\exists A \in \mathcal{X} \mid x \in A)$, therefore, $\bigcup \mathcal(X) = \mathbf(R)$

\end{document}